% The research corpus (ci/research.json): what a typical paper loads,
% in Times (mathptmx), Helvetica and Courier: the 35 base fonts.
\documentclass[11pt,a4paper]{article}
\usepackage[T1]{fontenc}
\usepackage{mathptmx}
\usepackage[scaled=0.92]{helvet}
\usepackage{courier}
\usepackage{amsmath}
\usepackage{graphicx}
\usepackage{booktabs}
\usepackage{siunitx}
\usepackage[ruled,linesnumbered]{algorithm2e}
\usepackage{listings}
\usepackage{xcolor}
\usepackage{tikz}
\usepackage[numbers,sort&compress]{natbib}
\usepackage{makeidx}
\usepackage{hyperref}
\makeindex
\lstset{basicstyle=\ttfamily\small,keywordstyle=\bfseries,columns=fullflexible,frame=single,numbers=left,numberstyle=\tiny}
\title{On the Typesetting of Research Papers in Times}
\author{A. Researcher\thanks{Department of Examples, University of Nowhere.} \and B. Colleague}
\date{1 January 2026}
\begin{document}
\maketitle
\begin{abstract}
We typeset a short paper with the fonts most journals ask for: Times for
the text and mathematics, Helvetica for headings in \textsf{sans serif},
and Courier for \texttt{code}. Tables, algorithms, listings, a picture,
citations and an index follow.
\end{abstract}

\section{Introduction}\label{sec:intro}
Literate programming\index{literate programming} \citep{knuth1984} and
document preparation\index{document preparation} \citep{lamport1994} are
older than most of their users. Structured programming
\citep{dijkstra1968} and information theory \citep{shannon1948,turing1938}
are older still. Section~\ref{sec:method} gives the method and
Table~\ref{tab:results} its results.

\section{Method}\label{sec:method}
Let $f\colon \mathbb{R}^n \to \mathbb{R}$ be convex\index{convexity}. Then
\begin{equation}
  f\Bigl(\sum_{i=1}^{k} \lambda_i x_i\Bigr) \le \sum_{i=1}^{k} \lambda_i f(x_i),
  \qquad \sum_{i=1}^k \lambda_i = 1,\ \lambda_i \ge 0 .
  \label{eq:jensen}
\end{equation}
Equation~\eqref{eq:jensen} is Jensen's inequality\index{Jensen's inequality}.
Algorithm~\ref{alg:gd} minimises $f$ by gradient descent\index{gradient descent}.

\begin{algorithm}
\caption{Gradient descent}\label{alg:gd}
\KwIn{$x_0$, step size $\eta > 0$, tolerance $\varepsilon$}
\KwOut{an approximate minimiser $x$}
$x \gets x_0$\;
\While{$\lVert \nabla f(x) \rVert > \varepsilon$}{
  $x \gets x - \eta \nabla f(x)$\;
}
\Return{$x$}\;
\end{algorithm}

\begin{lstlisting}[language=Python,caption={The same in Python}]
def descend(grad, x, eta=0.1, eps=1e-6):
    while norm(grad(x)) > eps:
        x = x - eta * grad(x)
    return x
\end{lstlisting}

\section{Results}
\begin{table}[ht]
\centering
\caption{Measured run times and speed-ups.}\label{tab:results}
\begin{tabular}{l S[table-format=3.2] S[table-format=1.2e1] r}
\toprule
Method & {Time (\si{\milli\second})} & {Error} & Iterations \\
\midrule
Baseline & 412.50 & 1.20e-3 & 1000 \\
Ours     &  87.25 & 3.40e-4 &  210 \\
\bottomrule
\end{tabular}
\end{table}
The speed-up is \num{4.73} at \SI{2.4}{\giga\hertz}; Figure~\ref{fig:pic}
draws the iterates.

\begin{figure}[ht]
\centering
\begin{tikzpicture}[scale=1.2]
  \draw[->] (-0.2,0) -- (4.2,0) node[right] {$x$};
  \draw[->] (0,-0.2) -- (0,3.2) node[above] {$f(x)$};
  \draw[thick,domain=0.2:3.9,smooth] plot (\x,{0.3*(\x-2)^2+0.5});
  \foreach \p in {0.4,1.1,1.6,1.9} \fill[red] (\p,{0.3*(\p-2)^2+0.5}) circle (1.5pt);
\end{tikzpicture}
\caption{Iterates of Algorithm~\ref{alg:gd} on a quadratic.}\label{fig:pic}
\end{figure}

\section{Conclusion}
\textbf{Bold}, \textit{italic}, \textsc{Small Caps}, \textsf{\textbf{bold
sans}} and \texttt{typewriter} all come from the base fonts. See
\url{https://www.ctan.org} for more.

\bibliographystyle{plainnat}
\bibliography{refs}
\printindex
\end{document}
